\subsection{Sprint 2}

Nesta sprint começamos com um impecilho, os AGES IV do time não possuíam conhecimento suficiente ou confiança para criar \textit{tasks} referente ao \textit{backend}, e como o 
único AGES III estava sobrecarregado no momento, fui solicitado a auxiliar nesse processo.
Trabalhei, inicialmente, na idealização dessas tasks, na implementação do endpoint GET para listar as aulas cadastradas no aplicativo e em seguida, 
desenvolvi o endpoint POST para criar comentários em aulas específicas.

Além do trabalho técnico na API, também me dediquei a atividades relacionadas à documentação da wiki do 
projeto. Realizei ajustes na página de \textit{home} e também criei a página referente ao banco de dados do projeto, contando com especificações e tecnologias utilizadas.
Por fim, devido à demanda e ao curto prazo, acabei revisando alguns merge requests (MRs) do frontend.

Em todas minhas atividades tive a companhia dos colegas AGES III Thomas Mello, AGES II Gabriel Suterio e AGES I Henrique Ribeiro.
\\
\\
\fbox{%
    \parbox{1\linewidth}{%
        \setlength{\parindent}{15pt}
            \itshape Devo dizer que pela minha parte, referente ao desenvolvimento da API, a visão que eu tinha do projeto era sempre positiva.
            Nada acabou sendo muito complexo ou trabalhoso, e pelo menos quem estava trabalhando junto comigo parecia estar gostando e aprendendo juntamente de mim.

            Porém, sinto que alguns fatores específicos contribuíram para o meu choque de realidade: nosso time tinha nove AGES I, apenas um AGES III (que tinha maior afinidade para o backend), e uma demanda muito maior para o frontend.
            Essa combinação mostrou ser extremamente desafiadora e estressante para meus colegas do frontend. Todos os débitos técnicos e atrasos que tivemos no projeto vieram desta área.

            Penso que uma decisão chave dos AGES IV no começo desta sprint também foi crítica para alguns problemas, especificamente de diversos colegas não contribuirem com o projeto.
            Não me recordo da explicação, se é que houve uma, mas nossos líderes desistiram da idea das squads e apenas seapraram o desenvolvimento do backend e frontend. A ideia era deixar
            cada membro escolher alguma task que quisesse trabalhar desde que fosse dentro da sua área.

            A partir daqui realmente houve uma decaída nas contribuições, e me colocando no lugar dos nossos AGES I sinto que era exatamente o problema que passei quando eu iniciei a AGES.
            Alguma divisão seja ela por pares, grupos ou squads sempre rende mais do que deixar os iniciantes por conta própria.
    }%
}
